\section {Introduction}

In the field of computer science education, understanding complex data structures and algorithms (DSA) often poses a significant challenge for learners. Traditional teaching methods, which rely heavily on static diagrams and theoretical explanations, can make it difficult for students to fully grasp dynamic processes such as pointer manipulation, hashing, tree traversal, and graph optimization. As a result, there is a growing need for more interactive and intuitive learning tools that can bridge the gap between theory and practice.

Motivated by this challenge, this project presents the design and development of a \textbf{DSA Visualizer}, an interactive application built using \textbf{C++} and the \textbf{SFML (Simple and Fast Multimedia Library)} framework. The primary goal of this project is to enhance the learning experience by providing real-time visual representations of fundamental data structures and algorithms. By allowing users to observe how operations are executed step by step, the system transforms abstract concepts into tangible and engaging visual experiences.

The visualizer focuses on four key areas: \textbf{Singly Linked Lists}, \textbf{Hash Tables}, \textbf{Tries (Prefix Trees)}, and \textbf{Minimum Spanning Trees (MST)}. These structures were carefully selected due to their importance in both academic curricula and real-world applications. Through interactive simulations, users can perform operations such as insertion, deletion, searching, and traversal, while simultaneously observing how the underlying data structure evolves.

This project is important because it not only improves conceptual understanding but also promotes active learning. By integrating a clear architectural separation between core logic and graphical rendering, the system ensures maintainability, scalability, and performance efficiency. Ultimately, the DSA Visualizer aims to serve as a practical educational tool that empowers students to explore, experiment, and develop a deeper intuition for algorithmic thinking.

\clearpage